\section{\arasim{} \label{app:AraSim}}

    Simulation of the ARA instrument, and Monte Carlo assessment of sensitivity, is achieved with a dedicated C++-based code called \arasim{}.

    The code is publicly available and hosted on GitHub.\footnote{\url{https://github.com/ara-software/AraSim}}
    Core dependencies of \arasim{} include the ROOT package~\cite{root}, 
    for class structure and data I/O,
    the Boost libraries, for pointer management, and FFTTools,\footnote{\url{https://github.com/nichol77/libRootFftwWrapper}}
    a FFTW-based package for managing fast Fourier transforms with ROOT wrappers.
    \arasim{} can also optionally link against \araroot{},\footnote{\url{https://github.com/ara-software/AraRoot}} the C++ software package for accessing ARA data and detector information.

    \arasim{} is designed to support end-to-end simulation of neutrino interactions in ARA.
    This includes the physics of the neutrino interaction (``event generation"),
    simulation of the Askaryan electric field, propagation of the field through the ice,
    receipt of the field by the antennas,
    application of the detector response and readout, including noise and trigger,
    and packaging of the simulated event in a data-like format.
        
    \subsection{Event Generation}
        
        \arasim{} supports three ways of generating events.
        First, \arasim{} contains a native Monte Carlo generator 
        of neutrino interactions.
        Second, users may provide a list of particle interactions and locations to be simulated.
        Third, users may directly provide the electric field at each antenna used in the simulation.
        In the first two cases, the Askaryan electric field from these interactions is generated as discussed in Appendix~\ref{app:askaryan_model}.
        
        \subsubsection{Native Generation of Events}

            \arasim{} was designed to be a backwards simulation of neutrinos.
            The user defines the radius and height of a fiducial cylinder, within 
            which neutrino interaction points are uniformly sampled.
            The neutrinos are then forced to interact, 
            warranting a weighting scheme to account for sampling bias.
            This unitless weight accounts for the probability that the neutrino 
            survives to its interaction point without an earlier interaction, and is defined as 
            \begin{equation}
              w_{\text{surv}} = \prod_i e^{-(\Delta D_i) / L_{\text{int,i}}}
              \label{eq:arasim_survival_weight}
            \end{equation}
            where the product is over each layer of the Earth model, $\Delta D_i$ is the distance a neutrino travels through that layer, and $L_{\text{int},i}$ is the interaction length of neutrinos in that layer.

            By default, the primary neutrino's flavor ($e$, $\mu$, or $\tau$), and whether it is a particle or anti-particle, is sampled uniformly over all six types of neutrinos.
            The interaction is modeled as a deep inelastic scattering
            via either neutral or charged currents, according to the energy-dependent branching ratios
            from~\cite{birefringence}.
            Users can change which neutrino flavors and interaction types are sampled for more specific studies.
            
            Several approximations are made to simplify the simulation of interaction processes.
            For example, \arasim{} simulates the decay of tau leptons only if the decay's energy deposition is greater than the cascade accompanying its production.
            In this case, only the signal generated by the decay is simulated.
            Additionally, Glashow resonance interactions are neglected by \arasim{}'s native event generator.
    
        \subsubsection{Event List Read-In\label{sec:eventlist}}

            A user may elect to simulate neutrino interactions with another framework but use \arasim{} to simulate the electric field from those interactions and the detector response.
            To do this, an event generator (\nuleptonsim{}, in this study~\cite{nuleptonsim}) is used to create a text file containing each interacting particle's flavor and whether it is an anti-particle; the interaction's type, inelasticity, location, and direction; and the primary neutrino's tracking ID and energy.
            \arasim{} uses the interaction's parameters to calculate the electric field from the interaction and processes the event through the
            remainder of the simulation steps (field propagation, detector simulation, etc.).
            If multiple events in the file list have the same primary neutrino tracking ID, these interactions are simulated as part of the same event.
            \arasim{} combines the signals from each of these interactions in a single voltage trace, delayed according to their relative interaction and propagation times.
            The resulting voltage trace is automatically lengthened to accommodate all of the contributing interactions.
        
        \subsubsection{Electric Field Read-In}

            A user may also provide a series of files to \arasim{} describing the Cartesian components of the electric field at each antenna as a function of time.
            This allows users to import results from other simulation frameworks, such as \texttt{FAERIE}~\cite{FAERIE}, \texttt{CoREAS}~\cite{COREAS}, and \texttt{Corsika~8}~\cite{corsika8} into \arasim{}.
            This allows a user to simulate the response of the
            ARA detector to non-neutrino physics signals,
            such as cosmic rays and beyond the Standard Model (BSM) processes, among others.

    \subsection{Askaryan Emission Model} \label{app:askaryan_model}

            \arasim{} supports two Askaryan emission models.
            They are described in detail elsewhere~\cite{hong_thesis},
            so we only briefly recount them here.
            The first is a ``time domain" model, where the electric field is found by time-differentiating the vector potential for the particle shower.
            The vector potential $\mathbf{A}(\theta, t)$
            under a far-field assumption has the following form:
            \begin{equation}
            \mathbf{A}(\theta, t) \propto\int_{-\infty}^{\infty} dz' Q(z') \cdot F_p\!\left(t, z'\right)
            \end{equation}
            where $z'$ is the shower depth in meters, $Q(z')$
            is the shower charge excess as a function of depth,
            $t$ is time, and $\theta$ is the angle between the shower axis
            and the RF signal propagation direction (viewing angle).
            $F_p$ is the ``shower form factor"~\cite{AZ-airshowers}, that 
            governs the shape of the radiated signal.
            The parameters of $F(p)$ are extracted from fits
            to the vector potentials from full simulations of 
            showers in media with ZHAireS~\cite{ARZ-signal}.
            Note that \arasim{} performs a full integral over the longitudinal
            profile of the shower, so that this model generates
            radiated signals with proper phase information.
            This ``time-domain" model, with correct phases,
            is the one used in this paper.

            The second Askaryan model is a ``frequency-domain" model,
            which is an implementation of the model in~\cite{showerenergy},
            and returns only the spectral amplitude as a function of frequency $f$:
            \begin{align}
              E(f) \propto \frac{f}{1150~\text{MHz}}\times \text{exp}\left[ - \frac{1}{2} \left( \frac{f}{1~\text{GHz}} \times \frac{\Omega}{2.2^{\circ}} \right)\right]
              \label{eq:alvarez-muniz}
            \end{align}
            The angular width of the cone of Askaryan radiation, $\Omega$, then follows the equations in Ref.~\cite{had-cone} for hadronic cascades and Ref~\cite{lpm_modeling} for electromagnetic cascades.
            Since this model only generates emission in the frequency
            domain, it does not account for the
            phases of the signal, and it therefore exaggerates the
            impulsiveness of the emission.
            This comparatively simple frequency-domain model is useful for comparisons to other simulation frameworks.

            In both emission models, the length of the electromagnetic particle shower is extended
            if the shower energy is greater than $10^{15}$~eV due to the Landau–Pomeranchuk–Migdal (LPM) effect following~\cite{lpm_modeling}.
            The LPM effect changes $F(p)$ in the time-domain model, while in the frequency-domain model it alters $\Omega$.
            
        
    \subsection{Modeling the Environment}

        After signals are generated, they must be propagated from the interaction vertex position to each of the station's channels.
        The properties of the Antarctic glacial ice have important effects on the propagation of radio waves, namely refractive, attenuation, and birefringence effects, all of which \arasim{} can simulate.
        Below, we describe how these effects are modeled, though birefringence effects~\cite{birefringence,ara_birefringence} are not considered in this study.
    
        \subsubsection{Index of Refraction}
            A refractive index model for the ice near ARA is required to determine the path of signal rays and their relative arrival times at each antenna.
            At leading order, the index of refraction $n$ depends only on depth from the surface $z$, which arises due to compactification
            of the glacial ice as new snow falls annually.
            The medium transitions from cold ($\sim -55~^\circ$C) loose snow at the surface, to relatively warm ($\sim -10~^\circ$C) ice at the bottom of the glacier~\cite{ice_temperature}.
            The implemented model for $n(z)$, which arises from ice-density arguments~\cite{ice_density}, takes the functional form
            of an exponential:

            \begin{align}
                n(z) = n_{\symrm{ice}} - (n_{\symrm{ice}}-n_{\symrm{firn}}) e^{-z/\Lambda}~,
            \end{align}
            %%
            where $n_{\symrm{ice}}$ and $n_{\symrm{firn}}$ are the ice's refractive indices in bulk and firn regions, respectively, and $\Lambda$ is the length scale of the transition between these two regions.
            Other functional forms for $n(z)$ are available in the literature (e.g.~\cite{ice_ara}), but are not currently supported.

            The values for $n_{\symrm{ice}}$, $n_{\symrm{firn}}$, and $\Lambda$
            must be derived from data.
            Several are available in \arasim{}, including those from the RICE experiment~\cite{ice_rice} and
            previous ARA studies~\cite{ara_a23, ara_pa_analysis}.
            For this paper, we adopted the model used in~\cite{ara_pa_analysis} with $n_{\symrm{ice}} = 1.780$, $n_{\symrm{firn}} = 1.326$, and $\Lambda = 0.0202^{-1}$~m~$\simeq 49.505$~m.

            Because the index of refraction depends on depth,
            Snell's law dictates that light paths follow curved, rather than rectilinear trajectories.
            This makes finding the paths non-trivial.
            To determine the ray paths, \arasim{} uses a numerical ray tracer that solves for the characteristics of the eikonal equation,
            \begin{align}
                \lvert \nabla T(\mathbf{r}) \rvert = n(\mathbf{r})~,
            \end{align}
            %%
            via a 4th-order Runge-Kutta integration, where $T$ is the propagation time, $n$ is the index of refraction, and $\mathbf{r}$ is a position in the ice.
            This numerical integration allows reflections at both the air-ice and ice-bedrock boundaries to be incorporated.
            Ray tracing is performed in two modes, depending on whether an initial-value or boundary-value problem is being solved.
            To solve an initial-value problem, an initial position and launch angle are provided, and the ray is determined via pure forward integration.
            More commonly, though, only an initial position (the interaction vertex location) and a final position (the antenna location) are provided.
            In this case, \arasim{} follows a two-step approach.
            The first ray-trace solution (the shortest, ``direct'', ray) is determined by calculating the launch angle via the semi-analytic method described in~\cite{ara_testbed}.
            Additional ray-trace solutions (in particular, the second-shortest, ``refracted'', ray) are found by determining other launch angles that minimize the distance of the forward-integrated ray to the target location via a root-finding algorithm.
            This hybrid approach is more computationally efficient compared to traditional trial-and-error methods.
            The ray tracer tracks the emission and receipt angles, 
            the path through the ice, and the travel time.
    
        \subsubsection{Ice Attenuation}
            Radio signals propagating in ice are attenuated over $\mathcal{O}(1)$~\text{km}~length scales~\cite{SPice-attenuation}.
            \arasim{} has several, increasingly detailed, implementations of the attenuation length.

            For this study, we used the most detailed implementation, which calculates the attenuation length as a function of both depth and frequency.
            The attenuation length's dependence on depth is due to the depth-dependence of the ice temperature, which in \arasim{} is parametrized via a cubic polynomial

            \begin{align}
                t(z) = az^3 + bz^2 + cz + d~,
            \end{align}
            %%
            where the $t$ is the temperature in degrees Celsius, $z$ is the distance from the ice surface, and the coefficients were determined from a fit to data from AMANDA and IceCube~\cite{ice_temperature}: $a=1.83415\times 10^{-9}$~$^\circ$C/m$^3$, $b=-1.59061\times 10^{-8}$~$^\circ$C/m$^2$, $c=2.67687\times 10^{-3}$~$^\circ$C/m, and $d=-51.0696$~$^\circ$C.

            Given an ice temperature, the attenuation length at three reference frequencies is modeled via a quadratic polynomial in temperature:

        \begin{align}\label{eq:reference_attenuation_lengths}
                \ln\left(\frac{1~\text{m}}{\lambda_{\symrm{att}}(f_i,t)}\right) = \alpha_i t^2 + \beta_i t + \gamma_i~,
            \end{align}
            %%
            where the coefficients, given in Table~\ref{tab:attenuation_length_coefficients}, were obtained by fitting data from~\cite{bogorodsky2012radioglaciology}.
            The final attenuation length is then obtained by linearly interpolating in log-frequency space between these reference frequencies:

            \begin{align}
                \ln\left(\frac{1~\text{m}}{\lambda_{\symrm{att}}(f)}\right) =~ &\ln\left(\frac{1~\text{m}}{\lambda_{\symrm{att}}(f_i)}\right) \nonumber \\ &+ \frac{\ln(f/f_i)}{\ln(f_{i+1}/f_i)} \ln\left(\frac{\lambda_{\symrm{att}}(f_i)}{\lambda_{\symrm{att}}(f_{i+1})}\right) ~,
            \end{align}
            %%
            where $i$ is such that $f_i \leq f < f_{i+1}$ and the dependence on temperature and depth has been suppressed for simplicity.

            \begin{table}[]
                \centering
                \begin{tabularx}{\linewidth}{c S S S}
                    \toprule
                    {$f_i$} & {$\alpha_i$ ($^\circ$C$^{-2}$)} & {$\beta_i$ ($^\circ$C$^{-1}$)} & {$\gamma_i$} \\
                    \midrule
                    $0.1$~MHz & -0.000884 & 0.026709 & -6.74890\\
                    $1$~GHz & -0.001773 & -0.070927 & -6.22121 \\
                    $3.16$~GHz & -0.000332 & -0.002213 & -4.09468 \\
                    \bottomrule
                \end{tabularx}
                \caption{Coefficients for the quadratic function~\eqref{eq:reference_attenuation_lengths} used to model the attenuation length at three reference energies: $f_0=0.1$~MHz, $f_1=1$~GHz, and $f_2=3.16$~GHz.
                \label{tab:attenuation_length_coefficients}}
            \end{table}

            To calculate the total attenuation factor experienced by a ray, \arasim{} calculates the product of attenuation factors in small steps along the ray path

            \begin{align}
                f_{\symrm{att}}(f) = \prod_i \exp\left(-\frac{L_i}{\lambda_{\symrm{att}}(f,z_i)}\right)~,
            \end{align}
            %%
            where $L_i$ is the length of the $i$th ray step and $z_i$ is the average depth in that step.
            This factor is then applied as $E_{\symrm{rec}}(f) = f_{\symrm{att}}(f) E_{\symrm{emit}}(f)$, where $E_{\symrm{rec}}$ is the frequency-space electric field amplitude arriving at the antenna and $E_{\symrm{emit}}$ is its amplitude at the emission point.

            \subsubsection{Birefringence}
            
            In addition, radio propagation in South Pole ice may be affected by anisotropies in the dielectric properties of the medium.
            In an anisotropic medium, the propagation speed depends on the local direction of propagation and the polarization of the signal.
            As a result, an incident electric field decomposes into two orthogonal propagation eigenstates, each of which experiences a different refractive index.

            The model implemented in \arasim{} follows the approach of~\cite{birefringence}, in which the anisotropy is described by three depth-dependent principal (refractive) indices, $(n_\alpha, n_\beta, n_\gamma)$, obtained from Ref.~\cite{ice_spicedata} and defined along orthogonal axes of the medium.
            In this framework, the local dielectric structure is represented by an indicatrix whose semi-axes are set by these three principal indices.
            For a signal propagating with local wave vector $\mathbf{k}$, the plane normal to $\mathbf{k}$ intersects the indicatrix in an ellipse.
            The principal axes of this ellipse define two local displacement-field eigenvectors, $\mathbf{D}_1$ and $\mathbf{D}_2$, and the corresponding semi-axis lengths determine the refractive indices $n_1$ and $n_2$ experienced by the two eigenstates.

            The ray path is determined by the standard ray tracing described above (i.e.\ the trajectory is calculated without birefringence corrections) and then subdivided into segments over which the propagation direction and local index of refraction are assumed to be constant.
            At each step, the local eigenstates and their associated refractive indices are calculated from the current propagation direction and the local principal indices.
            The accumulated time delay between the two eigenstates at the receiver is then

            \begin{align}
                \Delta T = \sum_i \frac{\Delta \ell_i}{c}\left(n_{2,i}-n_{1,i}\right),
            \end{align}
            %%
            where $\Delta \ell_i$ is the path length of segment $i$, and $n_{1,i}$ and $n_{2,i}$ are the effective refractive indices of the two eigenstates on that segment.

            Additionally, this model assumes that the local eigenstates evolve adiabatically along the ray path, so that the signal polarization continuously follows the rotation of the eigenstates.
            Therefore, birefringence modifies both the relative arrival times and the polarization of the received signals.
            
    \subsection{Modeling the Detector}\label{app:detector_modeling}
    
        After an electric field arrives at the station,
        \arasim{} must model the response of the detector.
    
        \subsubsection{Array and Station Layout}
        
            \arasim{} can be used to simulate up to $37$~identical ARA stations in a hexagonal layout with user-defined inter-station spacing.
            This capability supported the original design studies for the array.
            Simulation of ARA's Testbed, the five ``traditional" quad-string stations, is also supported.
            The Testbed is composed of 10 antennas deployed on 5 strings to a depth of at most 30~m~\cite{ara_design}.
            Typically, though, \arasim{} is used in a single-station simulation mode, where each station (A1--A5 and PA) is simulated using parameters imported from \araroot{}.
            These parameters include the station's calibrated antenna positions and other time-dependent properties, such as the exclusion of channels from the triggering logic.
            
            The PA detector, in particular, has three distinct simulation modes,
            representing the time evolution of the sub-detector.
            The first mode simulates the PA string as it was deployed in 2018,
            where it operated with seven VPol antennas and two HPol antennas.
            The second mode simulates the PA in 2019 when one VPol antenna from A5 was split and read out by both the A5 and PA DAQ boards.
            The third mode simulates the PA after 6 more VPol antennas from A5 were connected to the PA DAQ during the 2019--2020 season to aid in event reconstruction.
            
        \subsubsection{Antenna Gain Modeling \label{app:antenna}}

        After the electric field is propagated to the antenna location, \arasim{} converts the electric field into the voltage response of the antenna by convolving the signal with the antenna model in the frequency domain.

        For each ray solution, the electric field at the antenna is expressed in a local spherical basis $(\hat{r},\hat{\theta},\hat{\phi})$, where $\hat{r}$ points along the propagation direction of the incoming plane wave.
        VPol antennas are sensitive to the $\hat{\theta}$ component of the field, while HPol antennas are sensitive to the $\hat{\phi}$ component.
        
        The signal is then convolved with the antenna model frequency-by-frequency.
        For each frequency component $f$, the amplitude is scaled by the antenna vector effective height, while the phase is shifted according to the antenna phase response.
        In \arasim{}, the magnitude of the vector effective height is derived from the realized antenna gain according to
        \begin{equation}
        |\mathbf{h}_{\textrm eff}(f,\theta,\phi)|
        =
        \sqrt{
        \frac{G_r(f_{s},\theta,\phi)\, c^2\, Z_r}
        {4\pi\, n_{\textrm eff}\, f^2\, Z_0}
        }, \label{eq:heff}
        \end{equation}
        where $G_r(f_s,\theta,\phi)$ is the realized gain, $f_s = n_{\textrm eff} f / n$ is the medium-scaled frequency (see discussion below), $n_{\textrm eff}$ is the effective refractive index of the antenna in its borehole, $n$ is the refractive index in which the frequency-domain antenna gain was measured, $c$ is the speed of light in vacuum, $Z_0$ is the impedance of free space, and $Z_r$ is the receiver impedance.
        The vector effective height is the product of the magnitude with the
        relevant unit vector.
        For example, for a VPol antenna this would be $\mathbf{h}_{\textrm eff} = |\mathbf{h}_{\textrm eff}| \hat{\theta}$.
        The corresponding frequency-domain voltage at the antenna can then be written as
        \begin{equation}
        V_{\symrm{ant}}(f) = \mathbf{E}(f) \cdot \mathbf{h}_{\textrm eff}(f,\theta,\phi)\, e^{i\Phi_{\symrm{ant}}(f,\theta,\phi)},
        \end{equation}
        where $\mathbf{E}(f)$ is the incident electric-field and $\Phi_{\symrm{ant}}(f,\theta,\phi)$ is the antenna phase.

        For the gain of the antenna, \arasim{} supports several antenna-response model types, including models derived from electromagnetic simulations, from in-situ measurements, and anechoic chamber measurements.
        Separate responses are used for the bottom VPol, top VPol, and HPol antennas.
        For each antenna type, \arasim{} uses a realized gain
        $G_r(f,\theta,\phi)$
        and phase
        $\Phi_{\symrm{ant}}(f,\theta,\phi)$,
        where $f$ is the signal frequency and $(\theta,\phi)$ specifies the signal arrival direction in the antenna coordinate system.
        
        For this work, we use anechoic chamber measurements of ARA antenna responses taken at the University of Kansas~\cite{KU-antennas}.
        Since these responses were measured in air, while the detector operates in ice, an additional translation step is required before they can be applied in simulation.
        To account for this difference, \arasim{} translates the antenna response from the medium of the measurement to the local detector medium through a refractive-index scaling,
        \begin{equation}
        f_{\textrm s} = f\cdot \frac{n_{\textrm target}}{n_{\textrm source}},
        \end{equation}
        where $n_{\textrm source}$ is the refractive index of the medium in which the antenna response was measured, and $n_{\textrm target}$ is the refractive index of the medium in which the antenna is evaluated in simulation.
        This scaling accounts, at leading order, for the shift in the antenna's resonant frequency between media.
        
        The antenna model also includes information of impedance-mismatch of the antenna and its load through the standing-wave ratio (SWR), which is converted into a transmission coefficient and separately incorporated into the detector response:
        \begin{equation}
        \Gamma = \frac{\symrm{SWR}-1}{\symrm{SWR}+1},
        \qquad
        T = \sqrt{1-\Gamma^2},
        \end{equation}
        where $\Gamma$ is the voltage reflection coefficient and $T$ is the corresponding transmission factor.
        
        \subsubsection{Electronics Gain Modeling}

        After the antenna, \arasim{} models the response
        of the detector's amplifiers and filters.
        Similar to the antenna response, the detector electronics response is described through a frequency-dependent electronics gain $G_{\symrm{elec}}(f)$
        and phase
        $\phi_{\symrm{elec}}(f)$.
        \arasim{} supports either simulated electronics response models (e.g.\ \texttt{Qucs} simulations of the system) or in-situ, station- and configuration-dependent responses extracted from forced trigger data,\footnote{ARA stations use a supplementary $1$~Hz ``forced trigger'' to monitor the ambient background environment.} as shown in Figure~\ref{fig:gain_model}.

            \begin{figure}
                \centering
                \includegraphics[width=0.6\linewidth]{sections/araveff/images/veffPaper/gain_model_A3_C1.pdf}
                \caption{
                    \label{fig:gain_model}
                    Comparison between (linear) gain models of the system response in \arasim{}, as a function of frequency.
                    The previously used model, based on a \texttt{Qucs} simulation of the system (dotted), is compared to the data-driven models (solid lines, this work) used in this work, for the example of String 1 of A3 in configuration 1.
                    Data-driven models are presented for each channel: top VPol (TVPol), bottom VPol (BVPol), top HPol (THPol), and bottom HPol (BHPol).
                    }
            \end{figure}
            
        For each channel, the electronics response is applied in the frequency domain.
        The voltage spectrum before the electronics chain, $V_{\textrm ant}(f)$, is converted to the voltage spectrum after the electronics response as
        \begin{equation}
        V_{\textrm elec}(f) = V_{\textrm ant}(f)\, G_{\textrm elec}(f)\, e^{-i\phi_{\textrm elec}(f)},
        \end{equation}
        %%
        where $G_{\textrm elec}(f)$ and $\phi_{\textrm elec}(f)$ are the frequency-dependent electronics gain and phase response for the corresponding channel.

        To determine the in-situ gain model, forced trigger waveforms are Fourier transformed to obtain the measured voltage spectrum for each station, configuration, and channel.
        This measured spectrum, $H_{\symrm{meas}}(f)$, contains the thermal noise entering through the antenna convolved with the electronics gain.
        The expected thermal contribution at the antenna, $H_{\symrm{theory}}(f)$, is calculated from the antenna transmission coefficient and total noise temperature.
        The electronics gain model is then obtained by dividing the measured spectrum by the expected thermal spectrum,
        \begin{equation}
            G_{\symrm{elec}}(f) = \frac{H_{\symrm{meas}}(f)}{H_{\symrm{theory}}(f)}.
        \end{equation}
        %%
        \arasim{} recalculates the gain model to ensure it is consistent with the user-selected noise model (discussed in the following section) and antenna gain model.

        \subsubsection{Noise Modeling}

            After the antenna and electronics response have been applied, noise is added to the signal.
            In frequency space, thermal noise has spectral amplitudes drawn from a Rayleigh distribution.
            A noise model is therefore fully characterized by the scale parameter $\sigma(f)$ of the Rayleigh distribution at each frequency $f$.
            
            \arasim{} has two noise model implementations.
            The first is an idealized noise model,  
            which assumes white thermal noise across all frequencies.
            This is unrealistic, but it is useful for basic tests 
            of the simulation package or 
            to simplify assumptions for comparisons to other frameworks.
            In this idealized model, the Rayleigh distribution scale factor is determined assuming flat Johnson-Nyquist noise.
            For this purpose, we use a load resistance of $Z=50$~$\Omega$ and a total noise temperature of $T=325$~K to reflect both the $230$~K average ice temperature and the $95$~K noise temperature of the low noise amplifier (LNA).

            The second model, which is data-driven, is used in this study.
            To determine $\sigma(f)$ from data, the distribution of spectral amplitudes in frequency bins is generated from forced-trigger data.
            The resulting distribution is fit to a Rayleigh distribution to determine the scale factor, $\sigma_{r,a}(f)$, for each frequency bin and channel, $a$, on a station for that run, $r$.
            These values are averaged over all the runs in a livetime configuration to determine the final scale factor model $\sigma_{s,c,a}(f)$ for that channel $a$, livetime configuration $c$, and station $s$.
            
            Figure~\ref{fig:noise_model} shows a comparison of the noise models derived from in-situ data.
            These are shown after the system response has been applied.

            \arasim{} generates noise waveforms by sampling the frequency-space voltage amplitude from a Rayleigh distribution with scale $\sigma(f)$ given by the chosen model and phase sampled uniformly between $0$ and $2\pi$.
            The simulated noise trace is obtained by transforming the voltage back to the time-domain.
        
        \subsubsection{Traditional Trigger Simulation}

        Traditional ARA stations, A1--A5, use a
        multiplicity trigger.
        The trigger condition requires that impulsive power be observed independently in multiple channels
        within a coincidence window.
        In \arasim{}, the voltage waveform in each channel after the electronics response $V_{elec}(t)$ is first passed through a model of the tunnel diode, which acts as a power integrator.
        The diode output can be written as
        \begin{equation}
        D_i(t) = |V_{\symrm{elec}, i}(t)|^2 * h_{\textrm diode}(t),
        \end{equation}
        where $V_{\symrm{elect}, i}(t)$ denotes the voltage waveform in channel $i$, $h_{\textrm diode}(t)$ is the tunnel-diode response function, and $*$ denotes a convolution.
        A channel is considered to have triggered when the diode output exceeds a threshold $
        D_i(t) > D_{\textrm th}$.
        
        The detector triggers when a configurable number, typically three, of like-polarization channels satisfy this condition within a coincidence window $\Delta t_{\textrm trig}$.
        Physically, this reflects the expectation that a true impulsive plane-wave signal should appear across several antennas within a window, given roughly by the light-like travel time.
        The coincidence window is taken to be $\Delta t_{\textrm trig} \simeq 170~{\textrm ns}$, which corresponds to roughly the maximum time required for a wavefront to traverse the detector.
        
        \subsubsection{Phased Array Trigger Simulation}

        For the phased array station, the trigger is designed to improve sensitivity to low-amplitude signals using coherent beamforming.
        In a beamformed trigger, signals from $N$ closely spaced antennas are delayed and summed so that a plane wave arriving from a given direction adds coherently.
        For a coherent signal, the summed amplitude scales as $N$, while the RMS of incoherent thermal noise scales as $\sqrt{N}$, so that the $\symrm{SNR} \propto \sqrt{N}$.

        In \arasim{}, this trigger is represented through a trigger efficiency-based proxy rather than a full hardware-level simulation of beamforming.
        The simulation calculates the signal-only SNR in a short window in a reference VPol channel, chosen in the current implementation to be the top-most VPol antenna.

        This SNR is converted to the equivalent SNR for a signal arriving from the direction of the local calibration pulser, $\theta_{\symrm{CP}}$.
        This is done by scaling the measured SNR by the ratio of the PA's $50\%$ trigger efficiency SNR to signals from $\theta_{\symrm{CP}}$ to that from the signal arrival direction $\theta$,
        %%
        \begin{equation}
        {\textrm SNR}_{\textrm eff} = {\textrm SNR}_{\textrm meas}\, \frac{{\textrm SNR}_{50\%}(\theta_{\symrm{CP}})}{{\textrm SNR}_{50\%}(\theta)},
        \end{equation}
        %%
        where ${\textrm SNR}_{50\%}(\theta_{\symrm{CP}})=2$ and ${\textrm SNR}_{50\%}(\theta)$ is the angular-response described in~\cite{ara_pa_design}.
        The resulting effective SNR is mapped onto a trigger-efficiency curve $ \epsilon_{\textrm PA}({\textrm SNR}_{\textrm eff})$, determined from measurements.
        The PA then triggers randomly with a probability equal to $\epsilon_{\textrm PA}({\textrm SNR}_{\textrm eff})$.

        \subsubsection{Waveform Generation}

        \arasim{} saves a waveform for every successful trigger of the detector in simulation.
        Each trigger has an associated trigger time about which a readout window is defined.
        The size of the readout window, and how it is centered relative to the trigger time, is a configurable parameter that is set to match the settings of the real detector in the desired time period.
        The final voltage trace (after inclusion of signal, noise, and all detector effects) is then saved in the readout window for each channel.
        As discussed in Appendix~\ref{app:event_readout}, the saved simulated waveform is formatted to match real calibrated data.
    
    \subsection{Event Readout}\label{app:event_readout}

        \arasim{} saves all information using the ROOT file format.
        ROOT organizes files into ``trees,''
        where each tree contains ``branches'' that store individual variables or arrays.
        Each entry in a tree corresponds to a single event,
        allowing efficient, column-oriented access to large datasets
        without loading the entire file into memory.
        This structure makes it straightforward to selectively read
        only the branches relevant to a given analysis,
        which is particularly useful when processing large files.
        
        \subsubsection{Simulation Settings (\texttt{AraTree})}
        
        The simulation configuration is stored in the \texttt{AraTree}, which contains a single entry describing the detector configuration and other simulation conditions.
        This includes the detector geometry, the ice model, trigger configuration, and other simulation settings.
        In \arasim{}, \texttt{AraTree} is constructed with branches containing objects used to initialize the simulation:
        \begin{itemize}
            \item \texttt{Detector}, which defines the antenna layout and hardware responses,
            \item \texttt{IceModel}, which specifies the depth-dependent index of refraction and attenuation properties of the ice, 
            \item \texttt{Trigger}, which contains the trigger type and thresholds,
            \item \texttt{Settings}, which controls all simulation parameters,
            \item \texttt{Spectra}, which defines the neutrino flux model.
        \end{itemize}
        Since these objects are written once per simulation, \texttt{AraTree} serves as a record of the simulation configuration.
        
        \subsubsection{Simulation Results (\texttt{AraTree2})}
        
        The event-level simulation results are stored in \texttt{AraTree2}, which contains one entry per simulated neutrino event, or only triggered events upon request.
        Each entry includes both the generated event properties and the detector response.
        Specifically, \texttt{AraTree2} stores:
        \begin{itemize}
            \item The \texttt{Event} object, which carries the neutrino interaction properties such as vertex position, direction, energy, and interaction type,
            \item The \texttt{Report} object, which contains the detector response to that event.
        \end{itemize}

        The \texttt{Report} class holds information from all antennas and ray solutions.
        For each interaction and antenna, it stores quantities such as ray-tracing solutions between the vertex and antenna, propagation distances and attenuation, viewing and launch angles, polarization vectors, and the frequency-domain and time-domain voltage responses.
        If enabled, \arasim{} can store intermediate quantities, such as electric fields and frequency-domain signals prior to convolution with the detector response.
        
        \subsubsection{Results Mimicking Data (\texttt{eventTree})}
        
        In addition to simulation-specific output, \arasim{} produces a data-like representation of each event in \texttt{eventTree}.
        This tree is designed to mimic the structure of real ARA data and can be directly processed using analysis and reconstruction tools for ARA data.
        In particular, the \texttt{eventTree} contains a \texttt{UsefulAtriStationEvent}, which is the proprietary ARA data format,
        from \araroot{}, containing calibrated waveforms for each channel, 
        metadata, and an event weight corresponding to the simulated neutrino.
        These objects hold the final voltage waveforms after all detector effects have been applied, including signal propagation, antenna response, electronics response, and noise generation.

